% GENERATED FILE. Edit canonical lessons, then run npm run generate:books.
% canonical-source-hash: fnv1a64:0eebddcfbfe1d4dc
% canonical-lessons: MR-C446-vanyajiv, MR-C446-abhayaranya, MR-C446-bibatya, MR-C446-darkali-phodne, MR-C446-carne

\chapter{Wildlife Sanctuaries}
\label{ch:mr-wildlife-sanctuaries}
\hlchaptermodality{446}

\begin{chapteropening}
\textbf{By the end of this chapter:} \emph{I can talk about wildlife, a sanctuary and a leopard, and say a tiger roared and cattle grazed.}

Complete the last lesson of chapter 446: I can talk about wildlife, a sanctuary and a leopard, and say a tiger roared and cattle grazed.
\end{chapteropening}

\section[\texorpdfstring{\mr{वन्यजीव} (vanyajīv)}{वन्यजीव (vanyajīv)}]{\mr{वन्यजीव} (vanyajīv) — wildlife}

\label{lesson:MR-C446-vanyajiv}



\begin{quote}
Before the new one: say the Marathi for to blossom, then the Marathi for to wither.
\end{quote}

\subsection*{You'll want to know: \mr{वन्यजीव}}
\textbf{\mr{वन्यजीव}} — \emph{vanyajīv} — \textquotedblleft{}wildlife\textquotedblright{}.

\begin{quote}
\mr{या} \mr{जंगलात} \mr{खूप} \mr{वन्यजीव} \mr{आहेत}.
\end{quote}

\begin{quote}
\emph{yā jaṅgalāt khūp vanyajīv āhet.}
\end{quote}

\begin{quote}
There is a lot of wildlife in this forest.
\end{quote}

Masculine.

A formal word for the wild animals and birds of a region.

\begin{cousinweb}
\textbf{\mr{वन्यजीव}} (\emph{vanyajīv}) is a modern word built from Sanskrit \emph{vanya}, \textquotedblleft{}growing or living in a forest\textquotedblright{}, and \emph{jīva}, \textquotedblleft{}a living being\textquotedblright{} (Monier-Williams).
\end{cousinweb}

\begin{grammarlens}[title={What you've built}]
One more word for getting around.
\end{grammarlens}

\subsection*{Your turn}
\begin{itemize}
\raggedright
  \item \emph{Say it:} \emph{vanyajīv}
  \item \emph{Say it:} \emph{vanyajīv}, once more
  \item \emph{Say it:} \emph{yā jaṅgalāt khūp vanyajīv āhet.} — There is a lot of wildlife in this forest.
  \item \emph{Recall:} say the Marathi for to blossom, then the Marathi for to wither, then say \emph{vanyajīv} again
\end{itemize}

\begin{tcolorbox}[breakable,colback=teal!4,colframe=teal!35!black,title={Before you move on}]
What does \mr{वन्यजीव} mean? (\textquotedblleft{}Wildlife\textquotedblright{}.) Say it once more.
\end{tcolorbox}
\section[\texorpdfstring{\mr{अभयारण्य} (abhayāraṇya)}{अभयारण्य (abhayāraṇya)}]{\mr{अभयारण्य} (abhayāraṇya) — a sanctuary}

\label{lesson:MR-C446-abhayaranya}



\begin{quote}
Before the new one: say the Marathi for to wither, then the Marathi for wildlife.
\end{quote}

\subsection*{You'll want to know: \mr{अभयारण्य}}
\textbf{\mr{अभयारण्य}} — \emph{abhayāraṇya} — \textquotedblleft{}a sanctuary\textquotedblright{}.

\begin{quote}
\mr{अभयारण्यात} \mr{शिकार} \mr{करायला} \mr{बंदी} \mr{आहे}.
\end{quote}

\begin{quote}
\emph{abhayāraṇyāt śikār karāylā bandī āhe.}
\end{quote}

\begin{quote}
Hunting is banned in the sanctuary.
\end{quote}

Neuter.

\mr{अभय}, without fear, and \mr{अरण्य}, a forest.

\begin{cousinweb}
\textbf{\mr{अभयारण्य}} (\emph{abhayāraṇya}) is a modern word built from Sanskrit \emph{abhaya}, \textquotedblleft{}without fear, safe\textquotedblright{}, and \emph{araṇya}, \textquotedblleft{}a forest\textquotedblright{} (Monier-Williams): a forest where animals are safe.
\end{cousinweb}

\begin{grammarlens}[title={What you've built}]
One more word for getting around.
\end{grammarlens}

\subsection*{Your turn}
\begin{itemize}
\raggedright
  \item \emph{Say it:} \emph{abhayāraṇya}
  \item \emph{Say it:} \emph{abhayāraṇya}, once more
  \item \emph{Say it:} \emph{abhayāraṇyāt śikār karāylā bandī āhe.} — Hunting is banned in the sanctuary.
  \item \emph{Recall:} say the Marathi for to wither, then the Marathi for wildlife, then say \emph{abhayāraṇya} again
\end{itemize}

\begin{tcolorbox}[breakable,colback=teal!4,colframe=teal!35!black,title={Before you move on}]
What does \mr{अभयारण्य} mean? (\textquotedblleft{}A sanctuary\textquotedblright{}.) Say it once more.
\end{tcolorbox}
\section[\texorpdfstring{\mr{बिबट्या} (bibaṭyā)}{बिबट्या (bibaṭyā)}]{\mr{बिबट्या} (bibaṭyā) — a leopard}

\label{lesson:MR-C446-bibatya}



\begin{quote}
Before the new one: say the Marathi for wildlife, then the Marathi for a sanctuary.
\end{quote}

\subsection*{You'll want to know: \mr{बिबट्या}}
\textbf{\mr{बिबट्या}} — \emph{bibaṭyā} — \textquotedblleft{}a leopard\textquotedblright{}.

\begin{quote}
\mr{रात्री} \mr{गावाजवळ} \mr{बिबट्या} \mr{दिसला}.
\end{quote}

\begin{quote}
\emph{rātrī gāvājavaḷ bibaṭyā dislā.}
\end{quote}

\begin{quote}
A leopard was seen near the village at night.
\end{quote}

Masculine.

Before an ending \mr{बिबट्या} stays as it is: \mr{बिबट्याला}.

\begin{cousinweb}
\textbf{\mr{बिबट्या}} (\emph{bibaṭyā}) is the Marathi word for the leopard; its origin is not settled.
\end{cousinweb}

\begin{grammarlens}[title={What you've built}]
One more word for getting around.
\end{grammarlens}

\subsection*{Your turn}
\begin{itemize}
\raggedright
  \item \emph{Say it:} \emph{bibaṭyā}
  \item \emph{Say it:} \emph{bibaṭyā}, once more
  \item \emph{Say it:} \emph{rātrī gāvājavaḷ bibaṭyā dislā.} — A leopard was seen near the village at night.
  \item \emph{Recall:} say the Marathi for wildlife, then the Marathi for a sanctuary, then say \emph{bibaṭyā} again
\end{itemize}

\begin{tcolorbox}[breakable,colback=teal!4,colframe=teal!35!black,title={Before you move on}]
What does \mr{बिबट्या} mean? (\textquotedblleft{}A leopard\textquotedblright{}.) Say it once more.
\end{tcolorbox}
\section[\texorpdfstring{\mr{डरकाळी} \mr{फोडणे} (ḍarkāḷī phoḍṇe)}{डरकाळी फोडणे (ḍarkāḷī phoḍṇe)}]{\mr{डरकाळी} \mr{फोडणे} (ḍarkāḷī phoḍṇe) — to roar}

\label{lesson:MR-C446-darkali-phodne}



\begin{quote}
Before the new one: say the Marathi for a sanctuary, then the Marathi for a leopard.
\end{quote}

\subsection*{You'll want to know: \mr{डरकाळी} \mr{फोडणे}}
\textbf{\mr{डरकाळी} \mr{फोडणे}} — \emph{ḍarkāḷī phoḍṇe} — \textquotedblleft{}to roar\textquotedblright{}.

\begin{quote}
\mr{जंगलात} \mr{वाघाने} \mr{डरकाळी} \mr{फोडली}.
\end{quote}

\begin{quote}
\emph{jaṅgalāt vāghāne ḍarkāḷī phoḍlī.}
\end{quote}

\begin{quote}
In the forest the tiger roared.
\end{quote}

Past, with -\mr{ने} (he …): \mr{त्याने} \mr{डरकाळी} \mr{फोडली}, \emph{tyāne ḍarkāḷī phoḍlī}.

\mr{डरकाळी}, a roar, is feminine, so the past is \mr{डरकाळी} \mr{फोडली}.

\begin{cousinweb}
\textbf{\mr{डरकाळी} \mr{फोडणे}} (\emph{ḍarkāḷī phoḍṇe}) joins \textbf{\mr{डरकाळी}}, \textquotedblleft{}a roar\textquotedblright{}, whose origin is not settled, to \textbf{\mr{फोडणे}}, \textquotedblleft{}to break open\textquotedblright{}: to let out a roar.
\end{cousinweb}

\begin{grammarlens}[title={What you've built}]
One more word for getting around.
\end{grammarlens}

\subsection*{Your turn}
\begin{itemize}
\raggedright
  \item \emph{Say it:} \emph{ḍarkāḷī phoḍṇe}
  \item \emph{Say it:} \emph{ḍarkāḷī phoḍṇe}, once more
  \item \emph{Say it:} \emph{jaṅgalāt vāghāne ḍarkāḷī phoḍlī.} — In the forest the tiger roared.
  \item \emph{Recall:} say the Marathi for a sanctuary, then the Marathi for a leopard, then say \emph{ḍarkāḷī phoḍṇe} again
\end{itemize}

\begin{tcolorbox}[breakable,colback=teal!4,colframe=teal!35!black,title={Before you move on}]
What does \mr{डरकाळी} \mr{फोडणे} mean? (\textquotedblleft{}To roar\textquotedblright{}.) Say it once more.
\end{tcolorbox}
\section[\texorpdfstring{\mr{चरणे} (carṇe)}{चरणे (carṇe)}]{\mr{चरणे} (carṇe) — to graze}

\label{lesson:MR-C446-carne}



\begin{quote}
Before the new one: say the Marathi for a leopard, then the Marathi for to roar.
\end{quote}

\subsection*{You'll want to know: \mr{चरणे}}
\textbf{\mr{चरणे}} — \emph{carṇe} — \textquotedblleft{}to graze\textquotedblright{}.

\begin{quote}
\mr{डोंगरावर} \mr{गायी} \mr{चरत} \mr{होत्या}.
\end{quote}

\begin{quote}
\emph{ḍoṅgarāvar gāyī carat hotyā.}
\end{quote}

\begin{quote}
Cows were grazing on the hill.
\end{quote}

Past (he …): \mr{तो} \mr{चरला}, \emph{to carlā}.

Animals \mr{चरतात}; to take them grazing is \mr{चरायला} \mr{नेणे}.

\begin{cousinweb}
\textbf{\mr{चरणे}} (\emph{carṇe}) is inherited: CDIAL lists Marathi \emph{ċarṇẽ}, \textquotedblleft{}to move about, graze\textquotedblright{}, under Sanskrit \emph{carati}, \textquotedblleft{}moves, grazes\textquotedblright{} (CDIAL 4686).
\end{cousinweb}

\begin{grammarlens}[title={What you've built}]
One more word for getting around.
\end{grammarlens}

\subsection*{Your turn}
\begin{itemize}
\raggedright
  \item \emph{Say it:} \emph{carṇe}
  \item \emph{Say it:} \emph{carṇe}, once more
  \item \emph{Say it:} \emph{ḍoṅgarāvar gāyī carat hotyā.} — Cows were grazing on the hill.
  \item \emph{Recall:} say the Marathi for a leopard, then the Marathi for to roar, then say \emph{carṇe} again
\end{itemize}

\begin{tcolorbox}[breakable,colback=teal!4,colframe=teal!35!black,title={Before you move on}]
What does \mr{चरणे} mean? (\textquotedblleft{}To graze\textquotedblright{}.) Say it once more.
\end{tcolorbox}
